%! TeX root = ../main.tex

\sectionanchor{toc:photographs}{Photographs}

% ---------------------------------------------------------------------------
% Two pages. The first is a full-bleed frontispiece -- the studio portrait
% under a white scrim, the same device as the cover, with the type in the
% clear field at the foot. The second is an album page: three prints mounted
% with photo corners.
% ---------------------------------------------------------------------------

\pageghost{front-studio}

\vspace*{\fill}

\begin{center}
  {\Large\scshape\color{primary} Photographs}

  \vspace{3mm}
  {\color{rulecolor}\rule{18mm}{0.4pt}}

  \vspace{4mm}
  \begin{minipage}{0.82\linewidth}
    \centering\small\itshape\color{accent}
    ``What we have once enjoyed we can never lose; all that we love deeply
    becomes a part of us.''
  \end{minipage}
\end{center}

\vspace*{6mm}

\newpage

% Sizing is worked back from the 512pt text block: the two prints plus
% their captions and the gap come to roughly 471pt, leaving the page
% breathing rather than filled edge to edge -- which is what stops an album
% page looking like a contact sheet.
%
% No `center` environments here. Each one adds \topsep above and below, and
% two of them were enough (40pt) to push the lower pair onto a page of
% their own; \centering inside a group does the same job for nothing.
%
% The captions are not written here. \mountedphoto looks each one up by the
% plate it is handed, out of assets/data/captions.toml -- the file that
% captions every photograph in the booklet, wherever it is printed. To change
% what one of these prints says, edit its line there.
\vspace*{2mm}

{\centering
\mountedphoto[0.82\linewidth]{plates/plate-desk}\par}

\vspace{6mm}

% The two prints have different aspect ratios (1.378 and 1.411), so equal
% widths give unequal heights and the captions sit at different levels.
% Widths are set in that ratio instead -- 0.43 to 0.42 -- which brings both
% prints to the same 71mm height, so tops, bottoms and captions all line up.
{\centering
  \begin{minipage}[t]{0.43\linewidth}
    \centering
    \mountedphoto[\linewidth]{plates/plate-agbada}
  \end{minipage}%
  \hspace{0.04\linewidth}%
  \begin{minipage}[t]{0.42\linewidth}
    \centering
    \mountedphoto[\linewidth]{plates/plate-ledger}
\end{minipage}\par}

\vfill

\newpage

% The rest of his own photographs, out of assets/images/personal/, running
% straight on from the three prints above. No head over them and no break in
% the sequence: they are the same photographs of the same man, mounted the
% same way and set at about the same size, and the only reason these pages
% are generated and the one before was composed by hand is that there are
% too many of them to place one at a time.
%
% assets/build-personal.py sizes the prints and breaks the pages; drop a
% photograph of him into personal/, re-run assets/make-plates.sh, and these
% pages re-flow. Their captions come from assets/data/captions.toml, looked
% up by the file each photograph arrived in -- the same file the three above
% take theirs from.
% Generated by assets/build-personal.py -- do not edit by hand.
% Re-run assets/make-plates.sh after changing personal/.
%
% Provenance, so a photograph on the page can be traced back to
% the file it came from:
%   01  <-  WhatsApp Image 2026-08-31 at 16.59.18.jpeg
%   02  <-  WhatsApp Image 2026-08-31 at 17.01.57.jpeg
%   03  <-  WhatsApp Image 2026-08-31 at 21.13.20.jpeg
%   04  <-  album-04.jpg
%   05  <-  album-08.jpg
%   06  <-  album-32.jpg
%   07  <-  album-33.jpg
%   08  <-  franco_nero_knight_uniform.jpeg

\vspace*{8.13mm}
\personalrow{%
  \personalphoto{plates/personal/01}{52.50mm}%
  \personalgap
  \personalphoto{plates/personal/02}{52.47mm}%
}
\vspace{13.00mm}
\personalrow{%
  \personalphoto{plates/personal/03}{31.25mm}%
  \personalgap
  \personalphoto{plates/personal/04}{48.12mm}%
}

\newpage

\vspace*{15.50mm}
\personalrow{%
  \personalphoto{plates/personal/05}{44.81mm}%
  \personalgap
  \personalphoto{plates/personal/06}{61.98mm}%
}
\vspace{13.00mm}
\personalrow{%
  \personalphoto{plates/personal/07}{61.68mm}%
  \personalgap
  \personalphoto{plates/personal/08}{45.11mm}%
}



\vfill

\newpage

% The family's photographs, and the first head in the section: what it marks
% is the turn from his own pictures to everybody's. These carry a plain
% hairline frame rather than the photo corners used above -- the corners say
% "old print mounted in an album", which is true of his and not of these --
% and are set three across at half the size, because forty family group
% shots are a crowd and read as one.
\albumhead{Photo Album}

% Generated by assets/build-album.py -- do not edit by hand.
% Re-run assets/make-plates.sh after changing family_pics/.
%
% Provenance, so a photograph on the page can be traced back to
% the file it came from:
%   02  <-  WhatsApp Image 2026-08-31 at 16.58.30.jpeg
%   04  <-  WhatsApp Image 2026-09-08 at 09.01.34.jpeg
%   03  <-  WhatsApp Image 2026-08-31 at 17.06.14.jpeg
%   08  <-  album-05.jpg
%   09  <-  album-06.jpg
%   10  <-  album-07.jpg
%   01  <-  Screenshot_20220101-124628_WhatsApp.jpeg
%   05  <-  WhatsApp Image 2026-09-08 at 09.01.35.jpeg
%   06  <-  album-01.jpg
%   11  <-  album-09.jpg
%   13  <-  album-13.jpg
%   15  <-  album-16.jpg
%   07  <-  album-03.jpg
%   18  <-  album-19.jpg
%   17  <-  album-18.jpg
%   20  <-  album-21.jpg
%   19  <-  album-20.jpg
%   22  <-  album-28.jpg
%   12  <-  album-10.jpg
%   24  <-  album-31.jpg
%   21  <-  album-23.jpg
%   23  <-  album-30.jpg
%   25  <-  album-35.jpg
%   28  <-  franco_nero_and_emma.jpeg
%   29  <-  franco_nero_and_talia_n_tyshawn.jpeg
%   30  <-  franco_nero_and_treasure_n_princeton.jpeg
%   31  <-  franco_nero_helena.jpeg
%   27  <-  franco_nero_and_emike.jpeg
%   32  <-  franco_nero_with_sir_n_lady_chinwo.jpeg
%   14  <-  album-14.jpg
%   16  <-  album-17.jpg
%   26  <-  album-36.jpg

\vspace*{20.80mm}
\albumrow{%
  \albumphoto{plates/family/02}{38.61mm}{51.48mm}%
  \albumgap
  \albumphoto{plates/family/04}{77.87mm}{51.48mm}%
}
\vspace{7.00mm}
\albumrow{%
  \albumphoto{plates/family/03}{39.88mm}{53.17mm}%
  \albumgap
  \albumphoto{plates/family/08}{76.60mm}{53.17mm}%
}

\newpage

\albumrow{%
  \albumphoto{plates/family/09}{80.57mm}{52.97mm}%
  \albumgap
  \albumphoto{plates/family/10}{35.91mm}{52.97mm}%
}
\vspace{7.00mm}
\albumrow{%
  \albumphoto{plates/family/01}{66.25mm}{49.99mm}%
  \albumgap
  \albumphoto{plates/family/05}{50.23mm}{49.99mm}%
}
\vspace{7.00mm}
\albumrow{%
  \albumphoto{plates/family/06}{79.32mm}{55.23mm}%
  \albumgap
  \albumphoto{plates/family/11}{37.16mm}{55.23mm}%
}

\newpage

\vspace*{0.20mm}
\albumrow{%
  \albumphoto{plates/family/13}{77.98mm}{55.12mm}%
  \albumgap
  \albumphoto{plates/family/15}{38.50mm}{55.12mm}%
}
\vspace{7.00mm}
\albumrow{%
  \albumphoto{plates/family/07}{37.26mm}{55.65mm}%
  \albumgap
  \albumphoto{plates/family/18}{79.21mm}{55.65mm}%
}
\vspace{7.00mm}
\albumrow{%
  \albumphoto{plates/family/17}{61.49mm}{44.09mm}%
  \albumgap
  \albumphoto{plates/family/20}{54.98mm}{44.09mm}%
}

\newpage

\albumrow{%
  \albumphoto{plates/family/19}{77.35mm}{54.52mm}%
  \albumgap
  \albumphoto{plates/family/22}{39.12mm}{54.52mm}%
}
\vspace{7.00mm}
\albumrow{%
  \albumphoto{plates/family/12}{68.56mm}{49.02mm}%
  \albumgap
  \albumphoto{plates/family/24}{47.92mm}{49.02mm}%
}
\vspace{7.00mm}
\albumrow{%
  \albumphoto{plates/family/21}{37.34mm}{54.22mm}%
  \albumgap
  \albumphoto{plates/family/23}{37.67mm}{54.22mm}%
  \albumgap
  \albumphoto{plates/family/25}{38.25mm}{54.22mm}%
}

\newpage

\vspace*{29.27mm}
\albumrow{%
  \albumphoto{plates/family/28}{69.50mm}{51.75mm}%
  \albumgap
  \albumphoto{plates/family/29}{46.98mm}{51.75mm}%
}
\vspace{7.00mm}
\albumrow{%
  \albumphoto{plates/family/30}{38.49mm}{51.97mm}%
  \albumgap
  \albumphoto{plates/family/31}{77.99mm}{51.97mm}%
}

\newpage

\vspace*{41.87mm}
\albumrow{%
  \albumphoto{plates/family/27}{38.82mm}{51.75mm}%
  \albumgap
  \albumphoto{plates/family/32}{77.66mm}{51.75mm}%
}
\vspace{7.00mm}
\albumrow{%
  \albumphoto{plates/family/14}{37.73mm}{26.77mm}%
  \albumgap
  \albumphoto{plates/family/16}{37.46mm}{26.77mm}%
  \albumgap
  \albumphoto{plates/family/26}{38.07mm}{26.77mm}%
}



\vfill
